% GENERATED FILE. Edit canonical lessons, then run npm run generate:books.
% canonical-source-hash: fnv1a64:816c53d00c39f1fb
% canonical-lessons: RU-C507-rassada, RU-C507-tsvesti, RU-C507-vyanut, RU-C507-teplitsa, RU-C507-podrezat

\chapter{Seedlings and Blossom}
\label{ch:ru-seedlings-and-blossom}
\hlchaptermodality{507}

\begin{chapteropening}
\textbf{By the end of this chapter:} \emph{I can talk about growing flowers and trees.}

Complete the last lesson of chapter 507: I can talk about growing flowers and trees.
\end{chapteropening}

\section[\texorpdfstring{\ru{рассада} (rassáda)}{рассада (rassáda)}]{\ru{рассада} (rassáda) — seedlings}

\label{lesson:RU-C507-rassada}



\begin{quote}
Before the new one: say the Russian for a weed, then the Russian for a harvest.
\end{quote}

\subsection*{You'll want to know: \ru{рассада}}
\textbf{\ru{рассада}} (\emph{rassáda}) — \textquotedblleft{}seedlings\textquotedblright{}.

\begin{quote}
\ru{Весной} \ru{на} \ru{подоконнике} \ru{стоит} \ru{рассада} \ru{помидоров}.
\end{quote}

\begin{quote}
\emph{Vesnóy na podokónnike stoít rassáda pomidórov.}
\end{quote}

\begin{quote}
In spring there are tomato seedlings on the windowsill.
\end{quote}

Feminine, a collective noun.

\begin{cousinweb}
\emph{\ru{рассада}} is built on \emph{\ru{рассадить}} 'to plant out', from \emph{\ru{садить}} 'to plant'; EtymDB lists \emph{\ru{рассада}} without a source.
\end{cousinweb}

\begin{grammarlens}[title={What you've built}]
One more word for telling, explaining and describing at B1.
\end{grammarlens}

\subsection*{Your turn}
\begin{itemize}
\raggedright
  \item \emph{Say it:} \emph{rassáda}
  \item \emph{Say it:} \emph{rassáda}, once more
  \item \emph{Say it:} a sentence of your own with \emph{rassáda}
  \item \emph{Recall:} say the Russian for a weed, then the Russian for a harvest, then say \emph{rassáda} again
\end{itemize}

\begin{tcolorbox}[breakable,colback=teal!4,colframe=teal!35!black,title={Before you move on}]
What does \ru{рассада} mean? (\textquotedblleft{}Seedlings\textquotedblright{}.) Say it once more.
\end{tcolorbox}
\section[\texorpdfstring{\ru{цвести} (tsvestí)}{цвести (tsvestí)}]{\ru{цвести} (tsvestí) — to bloom}

\label{lesson:RU-C507-tsvesti}



\begin{quote}
Before the new one: say the Russian for a harvest, then the Russian for seedlings.
\end{quote}

\subsection*{You'll want to know: \ru{цвести}}
\textbf{\ru{цвести}} (\emph{tsvestí}) — \textquotedblleft{}to bloom\textquotedblright{}.

\begin{quote}
\ru{В} \ru{мае} \ru{в} \ru{нашем} \ru{саду} \ru{цветёт} \ru{вишня}.
\end{quote}

\begin{quote}
\emph{V máye v náshem sadú tsvetyót víshnya.}
\end{quote}

\begin{quote}
In May the cherry tree blossoms in our garden.
\end{quote}

Imperfective; present \ru{цветёт}, \ru{цветут} (\ru{ст} becomes \ru{т}). Perfective \ru{расцвести}.

\begin{cousinweb}
\emph{\ru{цвести}} is inherited from Proto-Slavic \emph{kvisti} (EtymDB), the root also behind \emph{\ru{цветок}} 'flower'.
\end{cousinweb}

\begin{grammarlens}[title={What you've built}]
One more word for telling, explaining and describing at B1.
\end{grammarlens}

\subsection*{Your turn}
\begin{itemize}
\raggedright
  \item \emph{Say it:} \emph{tsvestí}
  \item \emph{Say it:} \emph{tsvestí}, once more
  \item \emph{Say it:} a sentence of your own with \emph{tsvestí}
  \item \emph{Recall:} say the Russian for a harvest, then the Russian for seedlings, then say \emph{tsvestí} again
\end{itemize}

\begin{tcolorbox}[breakable,colback=teal!4,colframe=teal!35!black,title={Before you move on}]
What does \ru{цвести} mean? (\textquotedblleft{}To bloom\textquotedblright{}.) Say it once more.
\end{tcolorbox}
\section[\texorpdfstring{\ru{вянуть} (vyánut)}{вянуть (vyánut)}]{\ru{вянуть} (vyánut) — to wilt}

\label{lesson:RU-C507-vyanut}



\begin{quote}
Before the new one: say the Russian for seedlings, then the Russian for to bloom.
\end{quote}

\subsection*{You'll want to know: \ru{вянуть}}
\textbf{\ru{вянуть}} (\emph{vyánut}) — \textquotedblleft{}to wilt\textquotedblright{}.

\begin{quote}
\ru{Без} \ru{воды} \ru{цветы} \ru{быстро} \ru{вянут}.
\end{quote}

\begin{quote}
\emph{Bez vodý tsvetý býstro vyánut.}
\end{quote}

\begin{quote}
Without water flowers soon wilt.
\end{quote}

Imperfective; past \ru{вял} or \ru{вянул}. Perfective \ru{завянуть}.

\begin{cousinweb}
\emph{\ru{вянуть}} is inherited from Proto-Slavic \emph{vędnǫti} (EtymDB).
\end{cousinweb}

\begin{grammarlens}[title={What you've built}]
One more word for telling, explaining and describing at B1.
\end{grammarlens}

\subsection*{Your turn}
\begin{itemize}
\raggedright
  \item \emph{Say it:} \emph{vyánut}
  \item \emph{Say it:} \emph{vyánut}, once more
  \item \emph{Say it:} a sentence of your own with \emph{vyánut}
  \item \emph{Recall:} say the Russian for seedlings, then the Russian for to bloom, then say \emph{vyánut} again
\end{itemize}

\begin{tcolorbox}[breakable,colback=teal!4,colframe=teal!35!black,title={Before you move on}]
What does \ru{вянуть} mean? (\textquotedblleft{}To wilt\textquotedblright{}.) Say it once more.
\end{tcolorbox}
\section[\texorpdfstring{\ru{теплица} (teplítsa)}{теплица (teplítsa)}]{\ru{теплица} (teplítsa) — a greenhouse}

\label{lesson:RU-C507-teplitsa}



\begin{quote}
Before the new one: say the Russian for to bloom, then the Russian for to wilt.
\end{quote}

\subsection*{You'll want to know: \ru{теплица}}
\textbf{\ru{теплица}} (\emph{teplítsa}) — \textquotedblleft{}a greenhouse\textquotedblright{}.

\begin{quote}
\ru{В} \ru{теплице} \ru{у} \ru{нас} \ru{растут} \ru{огурцы}.
\end{quote}

\begin{quote}
\emph{V teplítse u nas rastút ogurtsý.}
\end{quote}

\begin{quote}
We grow cucumbers in the greenhouse.
\end{quote}

Feminine.

\begin{cousinweb}
\emph{\ru{теплица}} is built on the root of \emph{\ru{тёплый}} 'warm' with \emph{-\ru{ица}}; EtymDB lists the word without a source.
\end{cousinweb}

\begin{grammarlens}[title={What you've built}]
One more word for telling, explaining and describing at B1.
\end{grammarlens}

\subsection*{Your turn}
\begin{itemize}
\raggedright
  \item \emph{Say it:} \emph{teplítsa}
  \item \emph{Say it:} \emph{teplítsa}, once more
  \item \emph{Say it:} a sentence of your own with \emph{teplítsa}
  \item \emph{Recall:} say the Russian for to bloom, then the Russian for to wilt, then say \emph{teplítsa} again
\end{itemize}

\begin{tcolorbox}[breakable,colback=teal!4,colframe=teal!35!black,title={Before you move on}]
What does \ru{теплица} mean? (\textquotedblleft{}A greenhouse\textquotedblright{}.) Say it once more.
\end{tcolorbox}
\section[\texorpdfstring{\ru{подрезать} (podrezát)}{подрезать (podrezát)}]{\ru{подрезать} (podrezát) — to prune; to trim}

\label{lesson:RU-C507-podrezat}



\begin{quote}
Before the new one: say the Russian for to wilt, then the Russian for a greenhouse.
\end{quote}

\subsection*{You'll want to know: \ru{подрезать}}
\textbf{\ru{подрезать}} (\emph{podrezát}) — \textquotedblleft{}to prune; to trim\textquotedblright{}.

\begin{quote}
\ru{Осенью} \ru{дедушка} \ru{подрезает} \ru{старые} \ru{ветки}.
\end{quote}

\begin{quote}
\emph{Ósenyu dédushka podrezáyet stáryye vétki.}
\end{quote}

\begin{quote}
In the autumn Grandad prunes the old branches.
\end{quote}

Imperfective; perfective \ru{подрезать} (\ru{подрежу}), with the stress on the stem.

\begin{cousinweb}
\emph{\ru{подрезать}} is built on \emph{\ru{под}-} and \emph{\ru{резать}} 'to cut', which is inherited from Proto-Slavic \emph{rězati} (EtymDB).
\end{cousinweb}

\begin{grammarlens}[title={What you've built}]
One more word for telling, explaining and describing at B1.
\end{grammarlens}

\subsection*{Your turn}
\begin{itemize}
\raggedright
  \item \emph{Say it:} \emph{podrezát}
  \item \emph{Say it:} \emph{podrezát}, once more
  \item \emph{Say it:} a sentence of your own with \emph{podrezát}
  \item \emph{Recall:} say the Russian for to wilt, then the Russian for a greenhouse, then say \emph{podrezát} again
\end{itemize}

\begin{tcolorbox}[breakable,colback=teal!4,colframe=teal!35!black,title={Before you move on}]
What does \ru{подрезать} mean? (\textquotedblleft{}To prune; to trim\textquotedblright{}.) Say it once more.
\end{tcolorbox}
